\documentclass[a4paper,12pt]{article}
\usepackage[utf8]{inputenc}
\usepackage[bahasa]{babel} % Menggunakan Bahasa Indonesia
\usepackage{graphicx}
\usepackage{hyperref}
\usepackage{geometry}
\usepackage{array}
\usepackage{booktabs}
\usepackage{enumitem}

% Pengaturan Margin
\geometry{
    top=3cm,
    bottom=3cm,
    left=3cm,
    right=3cm
}

% Informasi Dokumen
\title{\textbf{Dokumentasi Kode Sumber dan Implementasi Sistem}}
\author{Lampiran Tesis - Temporal Graph Analytics}
\date{\today}

\begin{document}

\maketitle

\section*{Pengantar}
\label{sec:pengantar}

Dokumen ini berisi dokumentasi dan deskripsi fungsional dari seluruh artefak kode sumber (\textit{source code}) yang dikembangkan dalam penelitian tesis berjudul \textbf{"Temporal Graph Analytics for Cryptocurrency Portfolio Optimization"}.

Seluruh implementasi dilakukan menggunakan bahasa pemrograman Python 3.13 dalam lingkungan pengembangan Jupyter Notebook. Pendekatan ini dipilih untuk mendukung reprodusibilitas (\textit{reproducibility}) dan visualisasi data yang interaktif. Struktur kode dikategorikan menjadi modul utama, modul eksperimen lanjutan, dan eksperimen pendahulu.

\section{Modul Utama dan Perbandingan (Core Strategy)}
Bagian ini berisi kode inti yang digunakan untuk menghasilkan hasil akhir tesis, mencakup implementasi strategi \textit{Temporal Graph-Enhanced} dan perbandingannya dengan strategi \textit{Baseline}.

\begin{itemize}
    \item \textbf{\texttt{comparison\_baseline\_vs\_temporal\_graph\_fixed.ipynb}} \\
    \textit{(Final Product)} Notebook ini adalah \textbf{artefak utama} penelitian. Notebook ini mengintegrasikan seluruh komponen sistem:
    \begin{itemize}[label=$\circ$]
        \item Konstruksi \textit{Temporal Graph} dari data historis harga cryptocurrency.
        \item Ekstraksi fitur jaringan (\textit{Degree Centrality, Betweenness Centrality, Density}).
        \item Pelatihan model \textit{Machine Learning} (Random Forest) untuk deteksi rezim pasar.
        \item Simulasi \textit{backtesting} strategi Baseline vs Temporal Graph.
        \item Perhitungan metrik kinerja (Sharpe Ratio, Max Drawdown, Return) dan uji signifikansi.
    \end{itemize}
    
    \item \textbf{\texttt{temporal\_graph\_markowitz.ipynb}} \\
    Modul pengembangan ini berfokus pada validasi algoritma pembentukan graf dinamis. Digunakan untuk melakukan \textit{tuning} parameter ambang batas (\textit{threshold}) korelasi dan jendela waktu (\textit{window size}) sebelum diintegrasikan ke modul utama.
    
    \item \textbf{\texttt{ai\_gated\_markowitz\_complete.ipynb}} \\
    Implementasi mandiri dari strategi \textit{Baseline} (AI-Gated Markowitz Klasik). Berfungsi sebagai kontrol eksperimen untuk memastikan performa baseline terukur secara objektif tanpa bias dari modul graph.
\end{itemize}

\section{Eksperimen Pembanding Lanjutan (GNN)}
Eksperimen tambahan menggunakan \textit{Graph Neural Networks} (GNN) dilakukan untuk memvalidasi efektivitas pendekatan \textit{end-to-end deep learning} dibandingkan pendekatan \textit{feature extraction} manual yang diusulkan.

\begin{itemize}
    \item \textbf{\texttt{gnn\_enhanced\_comparison.ipynb}} \\
    Implementasi model \textit{Graph Convolutional Networks} (GCN) dan \textit{Graph Attention Networks} (GAT). Bertujuan membandingkan akurasi prediksi rezim pasar antara GNN dan Random Forest berbasis fitur graf.
    
    \item \textbf{\texttt{gnn\_simplified.ipynb}} \\
    Versi ringan dari implementasi GNN tanpa dependensi pustaka C++ yang kompleks, dikembangkan untuk validasi konsep cepat dan portabilitas kode.
\end{itemize}

\section{Eksperimen Pendahulu (Reinforcement Learning)}
Sebelum fokus pada pendekatan \textit{Temporal Graph}, penelitian mengeksplorasi metode \textit{Deep Reinforcement Learning} (DRL).

\begin{itemize}
    \item \textbf{\texttt{rl\_portfolio\_ultimate.ipynb}}: Implementasi agen PPO (\textit{Proximal Policy Optimization}) untuk manajemen portofolio. Disimpan sebagai dokumentasi \textit{negative results} atau metode pembanding awal.
\end{itemize}

\section{Ringkasan Berkas}
Daftar lengkap berkas kode sumber beserta deskripsi singkatnya disajikan dalam Tabel \ref{tab:daftar_notebook}.

\begin{table}[ht]
    \centering
    \caption{Ringkasan Artefak Kode Sumber Eksperimen}
    \label{tab:daftar_notebook}
    \renewcommand{\arraystretch}{1.5} % Menambah jarak vertikal antar baris
    \begin{tabular}{|p{6.5cm}|p{8.5cm}|}
    \hline
    \textbf{Nama Berkas Notebook} & \textbf{Deskripsi Fungsi Utama} \\ \hline
    
    \texttt{comparison\_baseline\_vs\newline\_temporal\_graph\_fixed.ipynb} & 
    \textbf{(FINAL)} Analisis komparatif lengkap antara Baseline dan Temporal Graph-Enhanced Strategy beserta visualisasi hasil. \\ \hline
    
    \texttt{temporal\_graph\newline\_markowitz.ipynb} & 
    Modul pengembangan algoritma konstruksi graf dinamis dan ekstraksi fitur sentralitas. \\ \hline
    
    \texttt{ai\_gated\_markowitz\newline\_complete.ipynb} & 
    Kode sumber strategi Baseline (Random Forest + Markowitz standar). \\ \hline
    
    \texttt{gnn\_enhanced\newline\_comparison.ipynb} & 
    Eksperimen pembanding menggunakan metode \textit{Graph Neural Networks} (GCN/GAT). \\ \hline
    
    \texttt{thesis\_experiment\newline\_final.ipynb} & 
    Notebook konsolidasi untuk agregasi data dan pelaporan metrik akhir. \\ \hline
    \end{tabular}
\end{table}

\end{document}
